\section{Introduction}
The neural separators that define the state of the art, such as
Demucs~\cite{defossez2020demucs}, Wave-U-Net~\cite{stoller2018waveunet}
and non-causal Conv-TasNet~\cite{luo2019convtasnet}, decide what a sample
is from a context that reaches into the future. A hearing aid or a hearable
cannot wait. Its output is heard together with the direct sound leaking
past the earpiece, and once the device delay grows beyond a few
milliseconds the listener hears comb filtering, echo, or a doubled own
voice. Tolerance studies put the usable budget at roughly 5--10\,ms in
total, and lower for open fittings~\cite{stone1999delay,stone2008opencanal}.
Such a device must therefore be \emph{strictly causal}, with a few
milliseconds of buffering at most. If it is binaural, it must also keep the
interaural time and level differences (ITD, ILD) that let the listener
place the talker and unmask speech in noise~\cite{doclo2015assisted}.

This paper prices that constraint. On \DataKind\ spatialised through
measured heads in simulated rooms, we sweep the \emph{algorithmic latency}
$L$ (synthesis window plus lookahead) from 1 to 32\,ms and measure, at each
$L$:
(i) how good \emph{any} mask can be (an oracle ratio mask, with symmetric
and asymmetric windows);
(ii) what causal networks reach: an asymmetric-window GRU with deep
filtering over past frames (AsymDF), the same network with a short
symmetric window, and a causal binaural Conv-TasNet;
(iii) what the same networks gain when allowed to hear the whole
utterance;
(iv) whether a fixed budget is better spent on window length or on
lookahead; and
(v) the measured CPU cost of streaming at batch size one.

Our hypothesis was that below about 8\,ms the loss is dominated by
frequency resolution rather than by the missing future. If so, spending
capacity on the \emph{past} (a long analysis window, recurrent state,
filtering over past frames) should recover most of it, keeping the 4\,ms
causal model within 1\,dB SI-SDRi of the 32\,ms one and its target cues
within the perceptual just-noticeable differences (JNDs). Section~\ref{sec:results}
reports which parts of this survive.
